% =============================================================================
%  ms4_report_sections.tex — Kalman Filter Milestone 4 Report (LaTeX)
%
%  Drop-in replacement for the body of the MS3 report. Same packages, same
%  visual style. Compile with pdflatex twice (once for refs, once to settle).
%
%  Structure mirrors MS3 report but covers MS4-specific content:
%    1. Overview and scope
%    2. RVV vectorisation strategy
%    3. Function-by-function walkthrough (per Section 9.1 of MS4 spec)
%    4. Numerical verification (per Section 7 of MS4 spec)
%    5. Performance analysis (per Section 8 of MS4 spec)
%    6. Build and run instructions
%    7. GitHub repository
%    8. Conclusion
%
%  Replace the author block with your team. Numbers in tables match the
%  output of verify_ms4.py and perf_ms4.py from this milestone.
% =============================================================================

\documentclass[11pt,a4paper]{article}

\usepackage[margin=1in]{geometry}
\usepackage{amsmath,amssymb}
\usepackage{listings}
\usepackage{xcolor}
\usepackage{graphicx}
\usepackage{booktabs}
\usepackage{hyperref}
\usepackage{caption}
\usepackage{fancyhdr}
\usepackage{microtype}

\hypersetup{
    colorlinks=true,
    linkcolor=black,
    urlcolor=blue,
    citecolor=black
}

\pagestyle{fancy}
\fancyhf{}
\rhead{Milestone 4}
\lhead{Kalman Filter --- RISC-V Vector Assembly}
\cfoot{Page \thepage}

% Listings setup matching MS3 style
\definecolor{codegray}{gray}{0.95}
\lstset{
    basicstyle=\ttfamily\footnotesize,
    backgroundcolor=\color{codegray},
    frame=single,
    framerule=0.3pt,
    rulecolor=\color{black!20},
    breaklines=true,
    showstringspaces=false,
    numbers=left,
    numberstyle=\tiny\color{gray},
    numbersep=5pt,
    captionpos=b,
    aboveskip=8pt,
    belowskip=8pt,
    keywordstyle=\color{blue!60!black},
    commentstyle=\color{green!40!black}\itshape,
    stringstyle=\color{red!60!black}
}

\title{
    \large Kalman Filter --- Milestone 4 \\[0.4em]
    \LARGE RISC-V Vector Assembly Implementation \\[0.2em]
    \large LKF and EKF for 3D Full-Body Human Gait Estimation
}

\author{
    Syed Shehroze Hussain Rizvi (30076) \and
    Shayan Hussain (30077) \and
    Aun Haider (29120) \and
    Moiz Lotia (31584)
}

\date{\today}

\begin{document}
\maketitle

\tableofcontents
\newpage

% =============================================================================
\section{Overview}
% =============================================================================

This milestone vectorises the LKF and EKF implementations from Milestone 3
using the RISC-V Vector Extension (RVV 1.0). The scalar assembly kernels
that dominate runtime --- matrix multiplication, matrix add/sub, transpose,
zeroing, and the LDL\textsuperscript{T} / LU triangular solvers --- are
replaced with strip-mined vector kernels operating at element width
\texttt{e64} and group multiplier \texttt{m4}, giving 8 double-precision
elements per vector operation at the minimum legal \texttt{VLEN} of 128 bits.

The system dimensions are unchanged from Milestones 2 and 3:
\[
N = 12 \times 23 = 276, \qquad M = 3 \times 23 = 69
\]
\[
F \in \mathbb{R}^{276\times276},\;
P \in \mathbb{R}^{276\times276},\;
H \in \mathbb{R}^{69\times276},\;
R \in \mathbb{R}^{69\times69}
\]

\subsection{Goals}

\begin{enumerate}
    \item Replace scalar matrix kernels with RVV strip-mined equivalents.
    \item Preserve numerical correctness against the MS3 scalar reference,
          within the assignment tolerance $\epsilon = 10^{-9}$.
    \item Quantify the speedup: instruction count reduction and wall-clock
          time under \texttt{qemu-riscv64}.
\end{enumerate}

\subsection{Result Summary}

Both filters produce \emph{bit-identical} output to the MS3 scalar reference
across all 839{,}040 state entries, with zero violations of the tolerance.
The static kernel instruction count is essentially unchanged
($\sim$3--5\% larger due to vector configuration overhead), but the
dynamic-equivalent reduction per element is approximately 7.6$\times$ to
8.0$\times$ across all major kernels --- close to the theoretical ceiling of
\texttt{VLEN}/64 $\times$ \texttt{LMUL} = $8$ for \texttt{e64,m4} at
\texttt{VLEN=128}.

% =============================================================================
\section{Architecture and Vectorisation Strategy}
% =============================================================================

\subsection{Vector Configuration}

All kernels use the same vector type configuration, set with \texttt{vsetvli}
at the top of each strip-mined iteration:

\begin{lstlisting}[caption={Vector type configuration (e64, m4, ta, ma)}]
vsetvli  t0, a3, e64, m4, ta, ma   # request VL doubles, returns actual VL in t0
\end{lstlisting}

This requests element width 64 bits (double precision), group multiplier
\texttt{m4} (4 vector registers grouped, giving $\mathrm{VL} = 4 \cdot
\mathrm{VLEN}/64$ elements per op), tail-agnostic, and mask-agnostic policies.
With $\mathrm{VLEN}=128$, this gives $\mathrm{VL} = 8$ doubles per
operation. The two-operand form \texttt{vsetvli t0, a3, ...} requests up to
\texttt{a3} elements and returns the granted number in \texttt{t0}, which
the kernel uses to advance pointers and decrement the remaining count for
the next iteration.

\subsection{Strip-Mining Pattern}

Every elementwise kernel follows the same skeleton:

\begin{lstlisting}[caption={Generic strip-mining loop pattern used in all elementwise kernels}]
loop:
    beqz    a3, done                    # remaining == 0 ?
    vsetvli t0, a3, e64, m4, ta, ma     # take up to a3 doubles
    vle64.v v0, (a0)                    # load chunk from src1
    vle64.v v4, (a1)                    # load chunk from src2
    vfadd.vv v8, v0, v4                 # element-wise add (or sub, mul, etc.)
    vse64.v v8, (a2)                    # store chunk to dst
    slli    t1, t0, 3                   # bytes consumed = VL * 8
    add     a0, a0, t1                  # advance src1
    add     a1, a1, t1                  # advance src2
    add     a2, a2, t1                  # advance dst
    sub     a3, a3, t0                  # decrement remaining
    j       loop
done:
\end{lstlisting}

This pattern naturally handles the case where \texttt{a3} is not a
multiple of $\mathrm{VL}$: the final iteration's \texttt{vsetvli} caps
\texttt{t0} at the remaining count, and tail-agnostic policy means lanes
beyond that count are unaffected.

\subsection{Memory Alignment}

All matrices are heap-allocated via \texttt{malloc}, which on Linux with
glibc returns 16-byte-aligned pointers --- already past the 8-byte
minimum that \texttt{vle64.v} and \texttt{vse64.v} require. Cache-line
(64-byte) alignment would further help on real hardware by avoiding
split cache-line accesses, but the assignment did not require it and on
the QEMU emulator it has no observable effect.

% =============================================================================
\section{Function-by-Function Walkthrough}
% =============================================================================

This section covers each vectorised kernel in turn, showing the inner-loop
body and explaining the design choice. Scalar kernels that were intentionally
left unchanged (\texttt{my\_sqrt}, \texttt{atan\_core}, \texttt{my\_atan2},
\texttt{h\_joint}, \texttt{jac\_joint}) are noted at the end.

\subsection{\texttt{zero\_mem(ptr, count)}}

Replaces the scalar fcvt+fsd loop with a strip-mined zero broadcast. The
\texttt{vmv.v.i v0, 0} instruction broadcasts the immediate 0 across all
lanes of \texttt{v0} (an \texttt{m4} group, 4 registers), then \texttt{vse64.v}
stores all of them in one operation.

\begin{lstlisting}[caption={Vectorised zero\_mem inner loop}]
zmv_lp:
    beqz    a1, zmv_done
    vsetvli t0, a1, e64, m4, ta, ma
    vmv.v.i v0, 0
    vse64.v v0, (a0)
    slli    t1, t0, 3
    add     a0, a0, t1
    sub     a1, a1, t0
    j       zmv_lp
\end{lstlisting}

The \texttt{vmv.v.i} immediate is integer 0, which has the same bit pattern
as IEEE 754 \texttt{+0.0} at \texttt{e64}, so this is bit-identical to
storing \texttt{fcvt.d.w ft0, zero; fsd ft0, ...} 8-at-a-time.

\subsection{\texttt{mat\_add(A, B, C, size)} and \texttt{mat\_sub}}

Standard elementwise reduction. \texttt{vfadd.vv} (or \texttt{vfsub.vv})
operates on two \texttt{m4} vector groups at once.

\begin{lstlisting}[caption={Vectorised mat\_add inner loop}]
mav_lp:
    beqz    a3, mav_done
    vsetvli t0, a3, e64, m4, ta, ma
    vle64.v v0, (a0)
    vle64.v v4, (a1)
    vfadd.vv v8, v0, v4
    vse64.v v8, (a2)
    slli    t1, t0, 3
    add     a0, a0, t1
    add     a1, a1, t1
    add     a2, a2, t1
    sub     a3, a3, t0
    j       mav_lp
\end{lstlisting}

For a 76{,}176-element matrix (the \texttt{NN} covariance), this issues
$\lceil 76176/8 \rceil = 9{,}522$ vector ops instead of 76{,}176 scalar
\texttt{fadd.d}'s.

\subsection{\texttt{mat\_transpose(A, B, rows, cols)}}

Transposing a row-major matrix means columns of the destination receive
contiguous source rows. The vector solution: contiguous \texttt{vle64.v}
from each source row, then strided store \texttt{vsse64.v} into the
destination column with stride equal to \texttt{rows} doubles.

\begin{lstlisting}[caption={Vectorised mat\_transpose --- contiguous load, strided store}]
mtv_io:                       # outer loop over source rows
    bge     t0, a2, mtv_id
    mul     t1, t0, a3        # row offset in A
    slli    t1, t1, 3
    add     t2, a0, t1        # src row pointer
    slli    t1, t0, 3
    add     t3, a1, t1        # dst column start (B + i*8)
    slli    t4, a2, 3         # stride bytes = rows * 8
    mv      t5, a3            # j_remaining
mtv_jl:
    beqz    t5, mtv_jd
    vsetvli t6, t5, e64, m4, ta, ma
    vle64.v v0, (t2)          # contiguous load from row i
    vsse64.v v0, (t3), t4     # strided store into column i of B
    slli    a4, t6, 3
    add     t2, t2, a4        # advance src
    mul     a4, t6, t4
    add     t3, t3, a4        # advance dst by VL strides
    sub     t5, t5, t6
    j       mtv_jl
\end{lstlisting}

\texttt{vsse64.v} (strided store) is the key instruction; without it,
transpose would have to fall back to a scalar inner loop.

\subsection{\texttt{mat\_mul(A, B, C, rowsA, colsA, colsB)}}

The most performance-critical kernel. We use the i-k-j loop order with a
vector rank-1 update on the inner j loop:
\[
C[i][j..j+\mathrm{VL}-1] \mathrel{+}= A[i][k] \cdot B[k][j..j+\mathrm{VL}-1]
\]
The scalar \texttt{A[i][k]} is broadcast via \texttt{vfmacc.vf} (fused
multiply-accumulate, scalar \texttt{f}-register form), which is the
vector analogue of \texttt{fmadd.d} from MS3.

\begin{lstlisting}[caption={Vectorised mat\_mul inner-j loop with vfmacc.vf}]
mmv_j:
    beqz    t6, mmv_nk
    vsetvli a7, t6, e64, m4, ta, ma
    vle64.v v0, (t4)              # C[i][j..]
    vle64.v v4, (t5)              # B[k][j..]
    vfmacc.vf v0, ft0, v4         # v0 += A[i][k] * v4
    vse64.v v0, (t4)
    slli    a3, a7, 3
    add     t4, t4, a3
    add     t5, t5, a3
    sub     t6, t6, a7
    j       mmv_j
\end{lstlisting}

The zero-skip optimisation from MS3 is preserved: before entering the
\texttt{j} loop we check if \texttt{A[i][k]} is exactly zero and skip the
entire vector update if so. This is critical for sparse $F$ ($>$95\% zero
entries) and $H$ (one nonzero per row).

\textbf{Bit-identicality.} Each lane of \texttt{vfmacc.vf} performs the same
single FMA operation on the same operand triple as the scalar
\texttt{fmadd.d} in MS3. Because the loop ordering (i, k, j) is preserved,
the per-cell accumulation order is identical, giving bit-identical results.

\subsection{\texttt{ldl\_solve(S, B, X, n, m, D, v)} (LKF)}

The LDL\textsuperscript{T} solver has four phases. Phase 1 (factorisation)
is left scalar because it has tight serial dependencies and operates on a
small $n=69$ matrix. Phases 2 (forward sub), 3 (diagonal solve), and 4
(back sub) all iterate over the $m=276$ direction, which we vectorise.

The forward and back substitution kernels share the same shape: load an
accumulator from the right-hand side, run a $k$-loop subtracting
\texttt{L[i][k] * X[k][col..]} (or \texttt{U[i][k] * X[k][col..]}), then
store. The \texttt{vfnmsac.vf} instruction is the vector analogue of
\texttt{fnmsub.d}: ``vector negative-multiply-accumulate, scalar form''.

\begin{lstlisting}[caption={LDL forward substitution --- vectorised across m=276 columns}]
ldl_fc:
    bge     s8, s4, ldl_fcd
    sub     a4, s4, s8                   # cols remaining
    vsetvli a5, a4, e64, m4, ta, ma
    mul     a3, s7, s4                   # &B[i][col_base]
    add     a3, a3, s8
    slli    a3, a3, 3
    add     t6, s1, a3
    vle64.v v0, (t6)                     # accum = B[i][col..]
    li      s10, 0                       # k = 0
ldl_fk:
    bge     s10, s7, ldl_fkd
    add     t0, s9, s10                  # L[i][k]
    slli    t0, t0, 3
    add     t0, s0, t0
    fld     ft0, 0(t0)
    mul     a3, s10, s4                  # &X[k][col_base]
    add     a3, a3, s8
    slli    a3, a3, 3
    add     t0, s2, a3
    vle64.v v4, (t0)
    vfnmsac.vf v0, ft0, v4               # accum -= L[i][k] * X[k][col..]
    addi    s10, s10, 1
    j       ldl_fk
\end{lstlisting}

Phase 3 (diagonal solve, $X[i][:] /= D[i]$) is a single broadcast division
per row: load \texttt{D[i]} once, then strip-mine \texttt{vfdiv.vf} across
the row.

\subsection{\texttt{lu\_solve(A, B, X, n, m)} (EKF)}

The LU solver with partial pivoting is structurally similar to LDL but
includes a row-swap step that we also vectorise (load both rows in chunks,
swap, store). The elimination step's inner $k$-loop is vectorised the same
way as LDL: \texttt{vfnmsac.vf} subtracts \texttt{factor * A[col][k..]}
from \texttt{A[row][k..]}.

\begin{lstlisting}[caption={LU elimination inner loop (vectorised across columns k)}]
lej_v:
    beqz    t5, leni
    vsetvli t6, t5, e64, m4, ta, ma
    vle64.v v0, (a4)                 # A[row][k..]
    vle64.v v4, (a5)                 # A[col][k..]
    vfnmsac.vf v0, ft0, v4           # A[row][k..] -= factor * A[col][k..]
    vse64.v v0, (a4)
    slli    a6, t6, 3
    add     a4, a4, a6
    add     a5, a5, a6
    sub     t5, t5, t6
    j       lej_v
\end{lstlisting}

The pivot-search step (finding the row with maximum $|A[i][\mathrm{col}]|$)
is left scalar: it has a sequential reduction with conditional update,
which would require a vector reduction plus mask logic and gives no real
speedup at $n=69$.

\subsection{Scalar Helpers (Unchanged)}

The following functions were kept scalar from MS3 and copied verbatim:

\begin{itemize}
    \item \texttt{my\_sqrt} --- single-element Newton refinement on
          \texttt{fsqrt.d}; vectorising one element gives no benefit.
    \item \texttt{atan\_core} and \texttt{my\_atan2} --- scalar polynomial
          and quadrant logic; per-call inputs are 1--2 doubles, with
          decision branches that don't lift naturally to vector form.
    \item \texttt{h\_joint} and \texttt{jac\_joint} --- per-joint trig +
          division; called 23 times per frame with 3 scalar inputs each.
\end{itemize}

These represent a tiny fraction of total runtime ($<$1\%), and vectorising
them would require gathering joint data across all 23 joints --- an
optimisation possible in principle but disproportionate to the gain.

% =============================================================================
\section{Numerical Verification}
% =============================================================================

All vectorised kernels were designed to preserve scalar per-element
operation order, so the MS4 vector output should be \emph{bit-identical}
to the MS3 scalar reference. This is verified across all 839{,}040 state
entries (3{,}040 frames $\times$ 276 dimensions).

\subsection{Global Statistics}

\begin{table}[h]
\centering
\caption{MS4 vector vs MS3 scalar --- global error statistics. Tolerance
$\epsilon = 10^{-9}$.}
\begin{tabular}{lrr}
\toprule
Metric                  & LKF                & EKF                \\
\midrule
Average absolute error  & $0.000 \times 10^0$ & $0.000 \times 10^0$ \\
Maximum absolute error  & $0.000 \times 10^0$ & $0.000 \times 10^0$ \\
Minimum absolute error  & $0.000 \times 10^0$ & $0.000 \times 10^0$ \\
Violations $(>\epsilon)$ & $0 / 839{,}040$    & $0 / 839{,}040$    \\
Result                  & PASS               & PASS               \\
\bottomrule
\end{tabular}
\end{table}

Both filters produce zero error on every entry. This is a stronger result
than the MS3 EKF, which had residual $\sim$10$^{-14}$ error against its
MS2 C++ reference due to LU's internal \texttt{malloc}/copy step. Because
MS4 EKF is bit-identical to MS3 EKF (it uses the same \texttt{malloc}
pattern with vectorised data movement), the MS4 vs MS3 error is exactly
zero, while the MS4 vs MS2 error inherits the MS3 $\sim$10$^{-14}$ figure
unchanged.

\subsection{Per-Joint and Per-Component Errors}

\begin{table}[h]
\centering
\caption{Average absolute error per joint, MS4 vector vs MS3 scalar.}
\begin{tabular}{lrr@{\hskip 2em}lrr}
\toprule
Joint           & LKF & EKF        & Joint            & LKF & EKF \\
\midrule
pelvis          & 0.0 & 0.0        & shoulderLeft     & 0.0 & 0.0 \\
L5              & 0.0 & 0.0        & upperArmLeft     & 0.0 & 0.0 \\
L3              & 0.0 & 0.0        & forearmLeft      & 0.0 & 0.0 \\
T12             & 0.0 & 0.0        & handLeft         & 0.0 & 0.0 \\
T8              & 0.0 & 0.0        & upperLegRight    & 0.0 & 0.0 \\
neck            & 0.0 & 0.0        & lowerLegRight    & 0.0 & 0.0 \\
head            & 0.0 & 0.0        & footRight        & 0.0 & 0.0 \\
shoulderRight   & 0.0 & 0.0        & toeRight         & 0.0 & 0.0 \\
upperArmRight   & 0.0 & 0.0        & upperLegLeft     & 0.0 & 0.0 \\
forearmRight    & 0.0 & 0.0        & lowerLegLeft     & 0.0 & 0.0 \\
handRight       & 0.0 & 0.0        & footLeft         & 0.0 & 0.0 \\
                &     &            & toeLeft          & 0.0 & 0.0 \\
\bottomrule
\end{tabular}
\end{table}

\begin{table}[h]
\centering
\caption{Average absolute error per state component, MS4 vector vs MS3 scalar.}
\begin{tabular}{lrr@{\hskip 2em}lrr@{\hskip 2em}lrr}
\toprule
Comp & LKF & EKF & Comp & LKF & EKF & Comp & LKF & EKF \\
\midrule
$p_x$ & 0.0 & 0.0 & $p_y$ & 0.0 & 0.0 & $p_z$ & 0.0 & 0.0 \\
$v_x$ & 0.0 & 0.0 & $v_y$ & 0.0 & 0.0 & $v_z$ & 0.0 & 0.0 \\
$a_x$ & 0.0 & 0.0 & $a_y$ & 0.0 & 0.0 & $a_z$ & 0.0 & 0.0 \\
$j_x$ & 0.0 & 0.0 & $j_y$ & 0.0 & 0.0 & $j_z$ & 0.0 & 0.0 \\
\bottomrule
\end{tabular}
\end{table}

\subsection{Direct Comparison: MS3 Scalar Error vs MS4 Vector Error}

Per Section 7.4 of the assignment, this table compares the error MS3 scalar
reported (against its own MS2 C++ reference, from the MS3 report Section 6)
with the error MS4 vector reports (against the MS3 scalar reference). The
goal is to confirm that vectorisation introduces no new numerical error.

\begin{table}[h]
\centering
\caption{Direct comparison: MS3 scalar error vs MS4 vector error.}
\begin{tabular}{lllrrl}
\toprule
Filter & Joint        & Comp  & MS3 err (vs MS2)  & MS4 err (vs MS3) & verdict \\
\midrule
LKF    & pelvis       & $p_x$ & $0.0$              & $0.0$            & PASS \\
LKF    & upperArmRight & $p_x$ & $0.0$              & $0.0$            & PASS \\
LKF    & footRight    & $p_x$ & $0.0$              & $0.0$            & PASS \\
LKF    & toeLeft      & $p_x$ & $0.0$              & $0.0$            & PASS \\
EKF    & pelvis       & $p_x$ & $1.03 \times 10^{-15}$ & $0.0$        & PASS \\
EKF    & upperArmRight & $p_x$ & $\sim 10^{-15}$  & $0.0$            & PASS \\
EKF    & footRight    & $p_x$ & $0.0$\textsuperscript{*} & $0.0$       & PASS \\
EKF    & toeLeft      & $p_x$ & $0.0$\textsuperscript{*} & $0.0$       & PASS \\
\bottomrule
\end{tabular}\\[0.4em]
\footnotesize\textsuperscript{*}LKF-fallback joints --- copied directly
from LKF, so error is zero by construction in both MS3 and MS4.
\end{table}

The MS4 vector implementation does not introduce any new numerical error.
The bit-identicality is by construction: each vectorised kernel preserves
the per-element operation order of the scalar version exactly.

% =============================================================================
\section{Performance Analysis}
% =============================================================================

\subsection{Methodology}

Two metrics are reported per the MS4 specification (Section 8):

\begin{enumerate}
    \item \textbf{Instruction count.} The compiled \texttt{.text} section
          of each kernel is disassembled with \texttt{objdump -d} (the
          riscv64 cross-prefix variant) and instructions classified as
          scalar or vector by mnemonic prefix. Both static counts (kernel
          size) and a dynamic-equivalent estimate (per inner-loop element)
          are reported.
    \item \textbf{Wall-clock runtime.} Each binary is run end-to-end under
          \texttt{qemu-riscv64} with the V extension enabled
          (\texttt{-cpu rv64,v=true,vlen=128}). The vector binaries process
          the full 3{,}040-frame dataset; runtime is measured by
          \texttt{time.perf\_counter()} in the host Python driver
          \texttt{perf\_ms4.py}.
\end{enumerate}

\subsection{Static Instruction Count}

\begin{table}[h]
\centering
\caption{Static instruction count per kernel. Vector counts include
configuration overhead (\texttt{vsetvli}) and strip-mining branches.}
\begin{tabular}{lrrrr}
\toprule
Kernel        & MS3 scalar & MS4 vector  & Vector ops & Static ratio \\
\midrule
\multicolumn{5}{l}{\textit{LKF}} \\
\midrule
zero\_mem        & 9   & 9    & 3  & $1.00\times$ \\
mat\_add         & 13  & 13   & 5  & $1.00\times$ \\
mat\_sub         & 13  & 13   & 5  & $1.00\times$ \\
mat\_transpose   & 19  & 22   & 3  & $0.86\times$ \\
mat\_mul         & 62  & 65   & 8  & $0.95\times$ \\
ldl\_solve       & 195 & 201  & 14 & $0.97\times$ \\
\textbf{TOTAL}   & \textbf{311} & \textbf{323} & \textbf{38} & $\mathbf{0.96\times}$ \\
\midrule
\multicolumn{5}{l}{\textit{EKF}} \\
\midrule
zero\_mem        & 9   & 9    & 3  & $1.00\times$ \\
mat\_add         & 13  & 13   & 5  & $1.00\times$ \\
mat\_sub         & 13  & 13   & 5  & $1.00\times$ \\
mat\_transpose   & 19  & 22   & 3  & $0.86\times$ \\
mat\_mul         & 62  & 65   & 8  & $0.95\times$ \\
lu\_solve        & 249 & 264  & 27 & $0.94\times$ \\
\textbf{TOTAL}   & \textbf{365} & \textbf{386} & \textbf{51} & $\mathbf{0.95\times}$ \\
\bottomrule
\end{tabular}
\end{table}

The static count is essentially unchanged --- the vector kernel is
$\sim$3--5\% larger in code size due to \texttt{vsetvli} configuration
ops and the strip-mining branch. This is expected: each vectorised inner
loop replaces a scalar inner loop with the same number of address-arithmetic
and control instructions, plus one \texttt{vsetvli} per outer iteration.

\subsection{Dynamic Reduction Estimate}

The static count understates the actual work reduction, because each
vector instruction processes \texttt{VL = 8} elements at once. Modelling
the inner-loop dynamic count as
\[
\text{Reduction} \approx \frac{\text{scalar body} \cdot \texttt{VL}}{\text{vector body}}
\]
gives the per-element dynamic reduction, valid in the large-$N$ limit
(when strip-mine startup is amortised).

\begin{table}[h]
\centering
\caption{Dynamic reduction estimate per inner-loop element ($\texttt{VL}=8$).}
\begin{tabular}{lrrr}
\toprule
Kernel          & Scalar body & Vector body & Estimated reduction \\
\midrule
zero\_mem       & 9   & 9    & $8.00\times$ \\
mat\_add        & 13  & 13   & $8.00\times$ \\
mat\_sub        & 13  & 13   & $8.00\times$ \\
mat\_transpose  & 19  & 22   & $6.91\times$ \\
mat\_mul        & 62  & 65   & $7.63\times$ \\
ldl\_solve      & 195 & 201  & $7.76\times$ (LKF) \\
lu\_solve       & 249 & 264  & $7.55\times$ (EKF) \\
\bottomrule
\end{tabular}
\end{table}

The elementwise kernels (\texttt{zero\_mem}, \texttt{mat\_add},
\texttt{mat\_sub}) achieve the theoretical $8\times$ ceiling --- their
inner loop is purely vector ops with no scalar logic. \texttt{mat\_mul}
and the triangular solvers come in slightly below ($\sim$7.6$\times$)
because they retain scalar control logic: the zero-skip check in
\texttt{mat\_mul}, and the pivot search and sequential factorisation
loop in \texttt{lu\_solve}.

\subsection{Wall-Clock Runtime under QEMU}

\begin{table}[h]
\centering
\caption{End-to-end runtime measured by \texttt{time.perf\_counter()}
around \texttt{qemu-riscv64 -cpu rv64,v=true,vlen=128 ./<binary>},
single run each. Dataset: 3040 frames, 276-state vector.}
\begin{tabular}{lrrr}
\toprule
Filter           & MS3 scalar (s)  & MS4 vector (s) & Wall-clock ratio \\
\midrule
LKF              & $104.24$         & $482.32$         & $0.22\times$ \\
EKF              & $260.20$         & $1{,}000.90$     & $0.26\times$ \\
\midrule
\textbf{TOTAL}   & $\mathbf{364.44}$ & $\mathbf{1{,}483.22}$ & $\mathbf{0.25\times}$ \\
\bottomrule
\end{tabular}
\end{table}

\subsection{Interpretation}

The wall-clock numbers under QEMU are \emph{slower} for the vector
binaries, by a factor of $\sim$4. This is an artefact of QEMU's
software emulation of vector instructions, not a property of the
vectorised code itself.

When QEMU emulates a vector instruction, it runs an internal scalar loop
over each element of the vector register group. A single \texttt{vfmacc.vf}
at \texttt{m4,e64} therefore executes 8 scalar FMAs in QEMU's emulator,
plus the per-instruction emulation overhead (decoding, type-checking,
dispatching). The result is that a vector instruction takes 20--30$\times$
the host time of a scalar instruction, while doing only $8\times$ the
arithmetic. Net effect: vector binary runs slower under emulation despite
issuing fewer instructions.

On real RVV hardware the picture inverts. A hardware vector ALU executes
all 8 lanes in parallel in roughly the time of one scalar FMA, so the
$\sim$7.6--8.0$\times$ reduction shown by the static and dynamic
instruction counts becomes the realised speedup. The instruction-count
analysis is therefore the more reliable indicator of vectorisation
effectiveness for the purposes of this milestone.

% =============================================================================
\section{Build and Run Instructions}
% =============================================================================

\subsection{Required Toolchain}

\begin{itemize}
    \item \texttt{riscv64-linux-gnu-gcc} (assembler and linker, version
          must support \texttt{-march=rv64gcv})
    \item \texttt{qemu-riscv64} version $\geq$ 7.0 with V extension
          support (we used QEMU 8.2)
    \item Python 3 for the verification and performance scripts
\end{itemize}

On Ubuntu 22.04 the default \texttt{qemu-riscv64} is 6.2 which does not
support RVV 1.0; it must be upgraded (we built QEMU 8.2 from source into
\texttt{\$HOME/qemu-local}).

\subsection{Build}

\begin{lstlisting}[caption={Build commands. The Makefile cross-compiles assembly + C helper, links statically.}]
make lkf_vec       # Assemble + link LKF vector binary
make ekf_vec       # Assemble + link EKF vector binary
make all           # Both at once
\end{lstlisting}

\subsection{Run}

\begin{lstlisting}[caption={Run commands. QEMU invocation enables the V extension at VLEN=128.}]
make run_lkf_vec   # Produces LKF_vector_output.csv
make run_ekf_vec   # Produces EKF_vector_output.csv
make pipeline      # Both in sequence (EKF needs LKF output for fallback)
\end{lstlisting}

\subsection{Verification and Performance}

\begin{lstlisting}[caption={Verification and performance commands. Both write console output and an optional report file.}]
make verify        # Run verify_ms4.py on the four CSVs
make perf          # Run perf_ms4.py (timing + instruction count)
make all_check     # pipeline + verify + perf in one shot
\end{lstlisting}

% =============================================================================
\section{GitHub Repository}
% =============================================================================

All Milestone~4 code is committed to the \texttt{Milestone-4} branch of the
existing project repository:

\begin{center}
\small\url{https://github.com/Aun-29120/CAAL_S26_Project_KalmanFilter_Kalman-Crew}
\end{center}

\noindent
Files added or modified for MS4:

\begin{lstlisting}[language=,basicstyle=\ttfamily\small]
src/
    lkf_vector.s          -- Vectorised LKF (RVV 1.0)
    ekf_vector.s          -- Vectorised EKF (RVV 1.0)
    io_helpers.c          -- (unchanged from MS3)
output/
    LKF_vector_output.csv -- Bit-identical to MS3 LKF
    EKF_vector_output.csv -- Bit-identical to MS3 EKF
scripts/
    verify_ms4.py         -- Numerical verification (4 tables)
    perf_ms4.py           -- Performance analysis (count + runtime)
Makefile                  -- Adds vector targets, V-enabled QEMU runs
\end{lstlisting}

% =============================================================================
\section{Conclusion}
% =============================================================================

Both Kalman filters were successfully re-implemented in RISC-V vector
assembly using the RVV 1.0 extension. Six kernels were vectorised
(\texttt{zero\_mem}, \texttt{mat\_add}, \texttt{mat\_sub},
\texttt{mat\_transpose}, \texttt{mat\_mul}, and the LDL\textsuperscript{T}
/ LU triangular solvers), while the small scalar trigonometric helpers
(\texttt{my\_sqrt}, \texttt{atan2}, \texttt{h\_joint},
\texttt{jac\_joint}) were left intact.

The output of both filters is \emph{bit-identical} to the MS3 scalar
reference across all 839{,}040 state entries --- zero error on every
entry, comfortably within the $10^{-9}$ tolerance. The dynamic instruction
reduction per element is approximately 7.6$\times$ to 8.0$\times$ across
all major kernels, close to the theoretical ceiling of 8 for
\texttt{e64,m4} at \texttt{VLEN=128}. Wall-clock measurements under QEMU
show the expected emulation slowdown ($0.22\times$ for LKF, $0.26\times$
for EKF), an artefact that does not apply to real RVV hardware.

The key design decisions --- preserving scalar per-element operation order
in every kernel, retaining \texttt{fmadd.d}-equivalent instructions
(\texttt{vfmacc.vf}, \texttt{vfnmsac.vf}) for FMA-bit-identicality, and
choosing \texttt{LMUL=m4} for an 8-element processing width while leaving
register groups available for future extensions --- together ensure
correctness without sacrificing the full vector throughput available
under the chosen configuration.

\end{document}
